%%%%%%%%%%  scan idx 236  |  folio 223  %%%%%%%%%%
est
bien déterminée (puisque c'est une sous-catégorie pleine de \(D(X)\)). Pour prouver l'existence, nous aurons besoin de quelques lemmes.

\textbf{Lemme 1.1.1.} \uline{Soit} \(i:X\to Y\) \uline{une immersion fermée de schémas telle que} \(i_*(\underline{O}_X)\) \uline{soit pseudo-cohérent}, \uline{et soit} \(E\in\ob D(X)\). \uline{Pour que} \(E\) \uline{soit \(n\)-pseudo-cohérent}, \uline{il faut et il suffit que} \(i_*(E)\) \uline{le soit}.

\uline{Preuve}. Le « il faut » résulte du critère général de stabilité de la pseudo-cohérence par application d'un foncteur (I~2.13). Inversement, supposons que \(i_*(E)\) est \(n\)-pseudo-cohérent, et prouvons qu'il en est de même de \(E\). Comme \(i_*(E)\) est à cohomologie localement bornée supérieurement, il en est de même de \(E\), donc, la question étant locale, on peut supposer que \(E\) est \(m\)-pseudo-cohérent pour \(m\) assez grand. Montrons, par récurrence descendante sur \(m\), que \(E\) est \(m\)-pseudo-cohérent pour tout \(m\geq n\). On peut supposer déjà prouvé que \(E\) est \((n+1)\)-pseudo-cohérent. Quitte à localiser, on a donc un triangle distingué
\[
\begin{tikzcd}
& E'' \arrow[dl,"|" marking] & \\
E' \arrow[rr] & & E \arrow[ul]
\end{tikzcd}
\]
où \(E'\) est parfait et \(E''^i=0\) pour \(i>n\). D'après (I~2.6), on est ramené à montrer que \(E''\) est \(n\)-pseudo-cohérent. Or, appliquant \(i_*\) au triangle précédent, on obtient un triangle distingué
\[
\begin{tikzcd}
& i_*(E'') \arrow[dl,"|" marking] & \\
i_*(E') \arrow[rr] & & i_*(E) \arrow[ul]
\end{tikzcd}\ .
\]
D'après la première partie de la démonstration, \(i_*(E')\) est pseudo-cohérent, donc en vertu de (I~2.6), \(i_*(E'')\) est \(n\)-pseudo-cohérent. Comme \(i_*(E'')\) est nul en degré \(>n\), cela signifie, d'après (I~2.10 b)), que \(H^n(i_*(E''))\) est de type fini. Comme \(H^n(i_*(E''))\simeq i_*H^n(E'')\), on est ramené à montrer, d'après (I~2.10 b)), que si \(M\) est un \(\underline{O}_X\)-Module tel que \(i_*(M)\) soit de type fini, alors \(M\) est de type fini, ce qui est immédiat (et indépendant de l'hypothèse sur \(i_*(\underline{O}_X)\)).

